% paper/sections/11f_interface_crossing.tex
%
% §11.5  不連続界面通過時の整合性
% ═══════════════════════════════════════════════════════════════════════════════
\subsection{不連続界面通過時の整合性}
\label{sec:interface_crossing}
% ═══════════════════════════════════════════════════════════════════════════════

密度比 $\rho_l/\rho_g = 1000$ の二相場において，
界面を跨ぐ速度場の連続性と圧力勾配の滑らかさを検証する．
GFM（§\ref{sec:gfm}）による鋭い界面処理が，
大きな物性不連続を正しくハンドリングすることの確認を目的とする．

% ─────────────────────────────────────────────────────────────────────────────
\subsubsection{変密度場における圧力の滑らかさ}
\label{sec:interface_crossing_detail}
% ─────────────────────────────────────────────────────────────────────────────

%% 実験結果は exp11_7_interface_crossing.py の実行後に挿入

\textbf{設定：}
計算領域 $[0,1]^2$，密度比 $\rho_l/\rho_g = 1000$，
界面位置 $x = 0.5$（平面界面）．
表面張力なし，重力なし（純粋な圧力 Poisson 問題のテスト）．
製造解法（MMS）により速度発散 $\bnabla\cdot\bu^*$ を処方し，
PPE を求解して得られる圧力場の界面近傍での挙動を観察する．
格子解像度 $N = 32, 64, 128$．

\textbf{評価指標：}
\begin{itemize}[noitemsep]
  \item 圧力 $p$ の界面法線方向プロファイル（$y = 0.5$ 上のラインプロット）
  \item $\bnabla p / \rho$ の界面近傍での連続性（速度補正項の滑らかさ）
  \item 格子収束次数
\end{itemize}

\textbf{合格基準：}
GFM 界面補正により，圧力勾配 $\bnabla p/\rho$ が界面で滑らかに接続され，
CSF（界面スミアリング）に比べて非物理的な振動が抑制されること．

%% TODO: exp11_7 結果（圧力プロファイル・収束テーブル）をここに挿入
